%%%%%%%%%%%%%%%%%%%%%%%%%%%%%%%%%%%%%%%%%%%%%%%%%%%%%%%%%%%%%%%%%%%%%%
% Lớp 8 | Chương II. Hằng đẳng thức đáng nhớ và ứng dụng | Phân tích đa thức thành nhân tử
% Dạng 8D2?9-3: Nhóm hạng tử, phối hợp nhiều phương pháp
%%%%%%%%%%%%%%%%%%%%%%%%%%%%%%%%%%%%%%%%%%%%%%%%%%%%%%%%%%%%%%%%%%%%%%


% nguồn: De_AI_Toan8_cuoi_hoc_ki_i_1 (nhập 21/11/2026)
\begin{bt}%[8D2V9-3]%[ma:2f975264]
	Phân tích các đa thức sau thành nhân tử:
a) $3x^3 - 12x$;
b) $x^2 + 6xy + 9y^2 - 16$;
c) $x^3 - 3x^2 - 4x + 12$.
	\loigiai{
		a) Đặt nhân tử chung $3x$ rồi dùng hằng đẳng thức hiệu hai bình phương: \\
$3x^3 - 12x = 3x(x^2 - 4) = 3x(x - 2)(x + 2)$. \\
b) Nhóm ba hạng tử đầu thành bình phương của một tổng rồi dùng hiệu hai bình phương: \\
$x^2 + 6xy + 9y^2 - 16 = (x + 3y)^2 - 4^2 = (x + 3y - 4)(x + 3y + 4)$. \\
c) Nhóm hạng tử rồi đặt nhân tử chung và dùng hằng đẳng thức: \\
$x^3 - 3x^2 - 4x + 12 = x^2(x - 3) - 4(x - 3) = (x - 3)(x^2 - 4) = (x - 3)(x - 2)(x + 2)$.
	}
\end{bt}
